% Section 7.3, from lectures/L07/appendix/b_exercises.md.

\section{Exercises}
\label{c7:sec:exercises}\label{c7:app:b}

\begin{aside}
\textbf{How to work these.} These exercises implement the \code{Uart} driver over the contracts of
\chapref{6}: the private register core that turns each access into one five-byte transaction, and
every public method of the interface built on it.

\textbf{How to check your work.} Nothing runs yet. The provided host suite,
\file{fw/test/uart/uart_test.cpp}, is guarded on four headers, two of which belong to \chapref{8},
so \code{make test} still reports there is nothing to do; the check on this work arrives in the next
chapter, and until then the specification below is the contract. \code{make build} has nothing to
build yet either, since it waits for the \file{source/main.cpp} that \chapref{8} provides, so check
that the driver compiles by hand, from \file{fw/}:
\code{g++ -std=c++17 -Wall -Wextra -Werror -Iinclude -fsyntax-only source/driver/uart/uart.cpp}.

Every exercise in this chapter is code, so Appendix~\ref{sol:ch} has nothing to add here.
\end{aside}

Work in \file{fw/}. You will add:

\begin{console}
include/
    driver/
        uart/
            uart.hpp
source/
    driver/
        uart/
            uart.cpp
\end{console}

The driver lives in the \code{driver::uart} namespace. These conventions apply to \file{include/}
and \file{source/}, which have to cross-compile for the ATmega; \file{test/} and
\file{include/arch/test/} are host-only and never reach avr-gcc, which is why the provided suites
next to them use \code{<cstdint>}, \code{std::}, \code{namespace driver::uart::test} and
\code{[[nodiscard]]} freely. Follow the same AVR-portable conventions as in \chapref{6}
(\secref{c6:sec:exercises}): \code{<stdint.h>} and bare \code{uint8_t} and \code{size_t} (no
\code{std::}), nested namespace blocks (not \code{namespace driver::uart { ... }}), and no
\code{[[nodiscard]]}.

% ------------------------------------------------------------------------------------------
\exerciseset{The \code{Uart} Driver}
\label{c7:set:1}

\exercise{The class}{C++}
\label{c7:ex:1.1}
In \file{include/driver/uart/uart.hpp}, declare \code{class Uart final : public Interface}. It holds
the transport as a reference member, \code{transport::Interface& myTransport}, injected through the
constructor, so that the same driver runs over the stub of \chapref{8} and over the real SPI of
\chapref{9}; the constructor therefore takes a \code{driver::transport::Interface&} and stores it.
Delete copy and move, since a reference member cannot be reassigned and the intent is worth making
explicit. Declare each \code{Interface} method with \code{override}, keeping the \code{noexcept} and
\code{const} qualifiers from the interface exactly.

Put the two private register helpers (Exercise~\ref{c7:ex:1.2}) in the private section. Implement
the bodies in \file{source/driver/uart/uart.cpp}.

\exercise{The register access core (private)}{C++}
\label{c7:ex:1.2}
Add two private helpers. They are the \textbf{only} methods that touch \code{myTransport}, and each
performs one five-byte SPI transaction from Part~3 of the protocol (\secref{proto:part3}).

\code{readReg(index)} returns a register's 32-bit value. It begins the transaction with
\code{myTransport.begin()}, then sends the command byte, which is the register index with the write
bit clear (bit 7 is 0), using \code{myTransport.transfer(command)} and ignoring the byte that comes
back, since the reply to the command byte is meaningless. It then clocks out the four data bytes by
calling \code{myTransport.transfer(0x00)} four times, sending dummy zeros, and assembles the four
returned bytes into the result \textbf{most significant byte first}, so the first byte back is bits
31 to 24. Finally it ends the transaction with \code{myTransport.end()} and returns the assembled
value.

Every byte here moves through the injected \code{myTransport}, using \code{myTransport.transfer()},
the raw duplex SPI exchange. Do \textbf{not} call the driver's own \code{write()}: that is the public
``send a UART byte'' operation, a different thing entirely, and calling it here would recurse.

\code{writeReg(index, value)} commits a 32-bit value to a register. It begins the transaction, sends
the command byte, which is the register index with the write bit set (bit 7 is 1), via
\code{myTransport.transfer(command)} while ignoring the reply, then sends the four bytes of
\code{value} \textbf{most significant first} (bits 31 to 24 first), each with
\code{myTransport.transfer(...)} and each reply ignored, and ends the transaction.

Make \code{readReg} a \code{const} member: it does not modify the \code{Uart} object; it exchanges
bytes through the transport, whose state is its own. \code{writeReg} need not be \code{const}.
Marking \code{readReg} \code{const} is exactly what lets \code{status()} and \code{errorFlags()} be
\code{const}.

Most significant first, in both helpers, is the detail the byte-order test exists to catch. A
reversed value is a different value to the hardware.

\task{Reaching a non-\code{const} transport from a \code{const} method (\code{const_cast})}
\code{readReg} is \code{const}, but the transport's \code{begin()}, \code{transfer()} and
\code{end()} are \textbf{not}: driving SPI is an action with side effects, not an observation. So a
\code{const} method that has to perform a register read must reach a non-\code{const} operation, and
this book writes that reach out explicitly:

\begin{cppcode}
uint32_t Uart::readReg(uint8_t addr) const noexcept
{
    auto& transport = const_cast<transport::Interface&>(myTransport);
    transport.begin();
    // ... transport.transfer(...) ...
    transport.end();
    // ...
}
\end{cppcode}

\noindent Read that cast carefully, because it is doing less than it looks like it is doing.
\textbf{The enclosing \code{const} never reached the transport in the first place.}
\code{myTransport} is a \emph{reference} member, and \code{const} does not propagate through a
reference: inside a \code{const} method, \code{myTransport} still names a plain
\code{transport::Interface&}, so the compiler would accept \code{myTransport.begin()} with no cast
at all. The cast yields the type the expression already had. In general, \code{const_cast} is what
you reach for when a logically \code{const} method (an observer such as \code{status()}) has to
invoke a non-\code{const} operation underneath, and you have decided the method should stay
\code{const} to its callers. It is a deliberate compromise, and worth understanding precisely.

The alternatives are marking the method non-\code{const}, which is honest but then means
\code{status()} cannot be \code{const}, or using a \code{mutable} member, which is right for a cache
and awkward for a reference. \code{const_cast} keeps the public \code{const} promise while admitting
that the implementation must drive something underneath. So why write a cast the compiler does not
ask for? Because it makes the intent explicit at the one line where a \code{const} method starts
driving hardware, and because it becomes genuinely \emph{required} the moment the transport is held
differently --- by value, or behind a \code{const}-propagating handle --- at which point the object's
\code{const} does reach the member and the calls stop compiling. Writing it now makes that change a
one-word edit rather than a redesign. Be clear which of the two you are relying on: here it is
documentation, not necessity.

The one hard rule is to cast away \code{const} only on an object that is not actually \code{const}.
\code{myTransport} refers to a real, non-\code{const} transport the caller owns, so driving it is
defined behaviour; casting away \code{const} on a genuinely \code{const} object and then mutating it
is undefined behaviour.

\exercise{The public methods}{C++}
\label{c7:ex:1.3}
Implement each \code{Interface} method in terms of the two helpers and the register map. The
register semantics of \chapref{5} make them short.

\code{configure(baudDiv)} sets the baud rate first and enables afterwards:
\code{writeReg(BAUD_DIV, baudDiv)}, then \code{writeReg(CTRL, 1U << ENABLE)}. Note the shift:
\code{ENABLE} is a bit \emph{position}, as every constant in the register map is, so writing it
directly would send \code{0} and clear the register rather than set bit 0.

\code{write(byte)} reads \code{STATUS}; if \code{TX_READY} is set it calls
\code{writeReg(TX_DATA, byte)} and returns \code{true}, and otherwise it writes nothing and returns
\code{false}.

\code{read(byte&)} reads \code{STATUS}; if \code{RX_VALID} is set it reads \code{RX_DATA} into the
out-parameter, then writes \code{writeReg(RX_POP, 1)} to advance the FIFO, and returns \code{true}.
If \code{RX_VALID} is clear it returns \code{false} and issues \textbf{no} \code{RX_POP}.

That success path is three distinct register accesses, in this order, which is the contract of
\secref{c5:app:a}. \textbf{Poll} \code{STATUS} to learn whether a byte is even available.
\textbf{Read} \code{RX_DATA}, which is a \emph{pure read}: it returns the front byte but does
\textbf{not} remove it from the FIFO. Then \textbf{pop} by writing \code{RX_POP}, a separate write
that discards the front byte so the next call sees the next one.

Because \code{RX_DATA} does not pop on its own, the \code{RX_POP} write is not optional. If you leave
it out, whether by forgetting it or by assuming the read already advanced the FIFO, the front byte is
never discarded, so every following \code{read} polls \code{RX_VALID}, finds it still set, re-reads
\code{RX_DATA}, and returns the \textbf{same byte forever}. Popping \emph{before} the read is the
opposite mistake: it discards the current byte unread and skips data.

The last three are one-liners. \code{status()} returns \code{readReg(STATUS)}, \code{errorFlags()}
returns \code{readReg(ERROR_FLAGS)}, and \code{clearErrors()} calls \code{writeReg(ERROR_FLAGS, 0)}.

\subsection*{What's ahead}
\chapref{8} scripts a transport under this driver: \code{driver::transport::Stub}, the provided host
suite that drives the real \code{Uart} over it, and the blocking helpers on top. Every mistake this
section warns about --- the reversed byte order, the missing \code{RX_POP}, the bit position used as
a mask --- is something a case in that suite exists to catch.
